\section{Minimal root covers for reduced boundaries}\label{cyc:section}

For a smooth projective pair of logarithmic Kodaira dimension zero
with reduced SNC boundary, the minimal cyclic cover has a
one-dimensional space of logarithmic pluriforms in every positive
degree. This strengthens
the character-line statement of Proposition~\ref{rankone:comparison}.
We prove the refinement together with its exact Hodge and divisorial
normalization. It gives a pure integral realization of the entire
highest line and fixes the normalization used in further
relative-Iitaka constructions.
The normalization lemmas below compare the minimum over actual fiber components
with the threshold over all divisorial valuations. They show that a
zero coefficient of the normalized base boundary yields a
multiplicity-one component, and that later permitted modifications
preserve the complete section spaces. These are the additional
normalization statements used by the variation companion.

\begin{lemma}[The cyclic cover of a logarithmic Kodaira-zero pair]
\label{lem:reduced-cyclic-top}
Let $(J,D)$ be a smooth geometrically integral projective pair over a
field $K$ of characteristic zero. Assume that $D$ is reduced and has
simple normal crossings and that
$\kappa(J_{\overline K},K_{J_{\overline K}}+D_{\overline K})=0$.
Choose the least positive integer $p$ for which
$H^0(J,p(K_J+D))\ne0$, a nonzero section $\omega$ in that space,
and put
\[
 E=\operatorname{div}(\omega)+pD,
 \qquad \Delta=D-\frac1pE.
\]
Write $\omega=b\eta^{\otimes p}$ in a rational canonical frame,
and let $\tau:N\to J$ be the normalization in
$K(J)[t]/(t^p-b)$. Let $\rho:\widetilde N\to N$ be a log
resolution, and set
\[
 P=\operatorname{Supp}(\Delta^{=1}),\qquad
 D_{\widetilde N}=\bigl((\tau\circ\rho)^{-1}P\bigr)_{\mathrm{red}}.
\]
Then $N$ is geometrically integral, the rational top form
$\Omega=t\tau^*\eta$ extends logarithmically to
$(\widetilde N,D_{\widetilde N})$, and
\[
 H^0\bigl(\widetilde N,m(K_{\widetilde N}+D_{\widetilde N})\bigr)
       =K\,\Omega^m\qquad(m\ge1).
\]
\end{lemma}

\begin{proof}
Every nonzero plurilogarithmic section space on $J$ is
one-dimensional: two independent sections of one degree would
define a nonconstant rational map. Consequently $\kappa(E)=0$,
and every effective integral divisor supported on $E$ also has
exactly one section, since it is bounded above by a multiple of $E$.
These assertions and the minimal index $p$ are unchanged by field
extension, by flat base change for global sections.

Over $\overline K(J)$, reducibility of $T^p-b$ would make $b$
an $\ell$th power for some prime $\ell\mid p$. The corresponding
root of $\omega$ would be a nonzero logarithmic pluriform of
degree $p/\ell$, since its logarithmic zero divisor is $E/\ell$.
This contradicts minimality. Thus $N$ is geometrically integral.

At a prime $Q$ of $N$ above a prime $R$ of $J$, write
\[
 a=\operatorname{coeff}_R E,\qquad
 d=\operatorname{coeff}_R D\in\{0,1\},\qquad
 e=\operatorname{ord}_Q(\tau^*R).
\]
The ramification formula for top differentials gives
\begin{equation}\label{eq:cyclic-root-order}
 \operatorname{ord}_Q\Omega=\frac{ea}{p}-ed+e-1.
\end{equation}
If $a=0$, the cover is unramified at $R$, and this order is
$-d$. If $a>0$, then $q=ea/p$ is a positive integer: the
integrality follows from $t^p=b$ and $d\in\mathbf Z$.
The order is $q-1\ge0$ when $d=1$, and $q+e-1\ge0$ when
$d=0$. Thus the divisorial poles are exactly the reduced inverse
image $D_N=(\tau^{-1}P)_{\mathrm{red}}$.

We also check extension across the exceptional divisors of $\rho$.
The subpair $(J,\Delta)$ is sub-log-canonical and is sub-klt
outside $P$. Indeed, its positive part is an SNC boundary bounded
by $D$, with all coefficients strictly below one outside $P$;
subtracting an effective divisor cannot worsen discrepancies.
Define the crepant finite pullback by
$K_N+\Delta_N=\tau^*(K_J+\Delta)$. With the compatible
canonical trivialization, $\Delta_N=-\operatorname{div}(\Omega)$.
For a prime divisor $V$ over $N$, let its restricted valuation
on $K(J)$ be $e_Vv$. The same ramification formula gives
\[
 1+\operatorname{ord}_V\Omega
       =e_V A_{J,\Delta}(v),
\]
where $A$ denotes log discrepancy. The right-hand side is
nonnegative, and is positive if the center lies outside $P$.
Since the left-hand side is integral, $\Omega$ has at most a
simple pole above $P$ and is regular elsewhere. This proves the
asserted logarithmic extension.

Finally let $H_N=\operatorname{div}_N(\Omega)+D_N$. Equation
\eqref{eq:cyclic-root-order} gives the sharper bound
\[
             0\le H_N\le\tau^*E.
\]
For completeness, when $a>0$ and $d=1$ its coefficient is
$q-1\le ea$; when $d=0$, it is
$q+e-1\le q+p-1\le pq=ea$, because $e\le p$ and $q\ge1$.
At $a=0$ its coefficient is zero.
Writing a logarithmic $m$-canonical form as $u\Omega^m$ and
testing its orders at the primes of the finite normalization gives
\[
 \operatorname{div}_N(u)+mH_N\ge0,
 \qquad
 u\in H^0(N,\mathcal O_N(m\tau^*E)).
\]
Finite pullback preserves Iitaka dimension, so
$\kappa(N,\tau^*E)=\kappa(J,E)=0$. Since $\tau^*E$ is
effective and $N$ is geometrically integral, the last space is
$K$. Conversely $\Omega^m$ is a logarithmic pluriform for every
$m\ge1$, proving the equality.
\end{proof}

\begin{corollary}[An integral pure realization of the moduli line]
\label{cor:reduced-integral-moduli}
Let $x:X\to W$ be a projective fibration between smooth complex
projective varieties, with geometrically integral generic fiber
$J$. Let $D_X$ be a reduced SNC boundary whose restriction $D_J$
satisfies the hypotheses of Lemma~\ref{lem:reduced-cyclic-top}.
Let $M_W$ be the rational Hodge divisor of the rank-one reduction,
with its rational frame and parabolic orders as in
Equation~\eqref{rankone:order}. On sufficiently high smooth models,
there is a polarizable integral pure variation
of Hodge structures on a dense smooth open $W^\circ\subset W$
whose top nonzero Hodge filtration piece is a line, and
\[
 M_W\sim_{\mathbf Q}
 \text{the rational canonical extension of that line}.
\]
In the canonical-bundle-formula normalization, this is the moduli
divisor of the auxiliary generic subpair
$\Delta_J=D_J-E_J/p$, which may have negative coefficients.
\end{corollary}

\begin{proof}
Let $s$ be the minimal-degree logarithmic generator on the generic
fiber, and choose a rational relative canonical frame $\eta$ on
$X$. The rational function $s/\eta^{\otimes p}$ defines one
cyclic cover over a dense open of $W$, not merely a collection
of local character covers. Spread its resolution and pole boundary
over a smaller dense open $W^\circ$ so that they form a smooth
projective family of SNC pairs. Write $r=\dim J$, and let
$u:\widetilde X^\circ\setminus D_{\widetilde X}^\circ\to W^\circ$
be the complementary open family. Its geometric variation
\[
 \mathbb V_{\mathbf Z}=R^ru_*\mathbf Z/\mathrm{torsion}
\]
is admissible and graded-polarizable. Logarithmic Hodge theory
identifies
\[
 F^{r+1}\mathbb V=0,\qquad
 F^r\mathbb V_{\mathcal O}
   =\widetilde x_*\omega_{\widetilde X^\circ/W^\circ}
                           (D_{\widetilde X}^\circ).
\]
Lemma~\ref{lem:reduced-cyclic-top} makes this entire top piece a
line. Strictness for the weight filtration therefore singles out
one weight $w$ with
\[
 F^rW_{w-1}\mathbb V=0,
 \qquad F^rW_w\mathbb V=F^r\mathbb V,
 \qquad
 F^r\operatorname{Gr}^W_w\mathbb V\simeq F^r\mathbb V.
\]
The rational variation $\operatorname{Gr}^W_w\mathbb V$ is
polarizable. Intersecting its rational weight filtration with
$\mathbb V_{\mathbf Z}$ and removing torsion supplies an integral
lattice; after rescaling, its rational polarization is integral.
Thus this pure variation has precisely the required top line.

After a finite cover making local monodromies unipotent, the
canonical extensions have locally free double gradeds for the
Hodge and weight filtrations. The displayed identification of top
lines consequently extends across the boundary; see
\cite[Proposition~3.12]{FujinoFujisawa}. Descent gives the same
identification for their rational canonical extensions.

The global rational frame of this entire top line is
$s^{1/p}$. Equation~\eqref{rankone:order} computes its parabolic
order at every prime divisor $P$ as $\beta_P$, the order of the
same rational frame in $M_W$. Equality of these divisorial
orders identifies the rational extensions. In particular, it
identifies $M_W$ itself, with no additional factor of $p$.
We check the rank condition for the canonical-bundle-formula
comparison. On a log resolution $\mu:J'\to J$, let
$R'=\lceil A^*(J,\Delta_J)_{J'}\rceil$, where $A^*$ omits the
log-discrepancy-zero components of the discrepancy divisor.
The subpair is sub-log-canonical, so $R'\geq0$. At the strict
transform of each original prime its coefficient is zero if
$0\leq\Delta_J\leq1$, and is $\lceil-\Delta_J\rceil$ otherwise.
Thus $0\leq\mu_*R'\leq E_J$. Testing original primes gives
$H^0(J',\cO_{J'}(R'))\subseteq H^0(J,\cO_J(E_J))=\C(W)$,
and constants give equality. This is the required rank-one condition.
Extending the generic subpair by
$-p^{-1}\operatorname{div}_{K_{X/W}}(s)$ makes its log canonical
divisor rationally equivalent to $x^*K_W$. The generic sub-lc,
relative triviality, and rank conditions of
\cite[Definition~6.18]{BFMT} therefore hold. Adding a base pullback
to normalize vertical orders changes the total base term and the
discriminant equally, leaving the moduli divisor unchanged.

The least positive index trivializing $K_J+\Delta_J$ is $p$.
Indeed, a smaller such index $q$ would produce a rational
$q$-canonical form with logarithmic zero divisor $qE_J/p\geq0$,
contradicting the minimality of $p$. Thus the cyclic field extension
in Construction-Definition~6.23 of \cite{BFMT} is the one used above.

For its boundary comparison, put
$P_J=\operatorname{Supp}(\Delta_J^{=1})$ and first make a common
crepant log resolution of the auxiliary subpair and its root cover. The
coefficient-one support on the crepant model is contained in the
total inverse image of $P_J$: outside $P_J$ the auxiliary subpair is
sub-klt. Every additional component allowed by our pole boundary
has positive log discrepancy, so the ramification discrepancy
formula in Lemma~\ref{lem:reduced-cyclic-top} makes $\Omega$
regular there. The top-form space for the BFMT pole boundary
therefore injects into the top-form space for our boundary and
contains $\Omega$; both are the line generated by $\Omega$.
After spreading over a common smooth base open, inclusion of the
open complements induces a morphism of their admissible geometric
mixed variations. This morphism is strict for the Hodge and weight
filtrations, including their limiting filtrations after a finite
cover making boundary monodromy unipotent. It therefore identifies
the top lines in the same pure weight grade. Functoriality
of canonical extension gives the same identification across the
boundary and hence for the parabolic rational lines. Theorem~6.28 of
\cite{BFMT} now identifies this Hodge divisor with the moduli
divisor using one common positive multiple on both sides. This
comparison permits negative coefficients in the generic subpair;
it does not use the separate semiampleness theorem requiring
generic effectiveness.
\end{proof}

\begin{remark}
The hypothesis that $D$ is reduced is essential to
Lemma~\ref{lem:reduced-cyclic-top}. For five distinct points on
$\mathbf P^1$, the rational boundary
$D=\frac25(P_1+\cdots+P_5)$ satisfies $K_{\mathbf P^1}+D\sim_{\mathbf Q}0$.
Its index-five root cover has genus six and empty pole boundary.
Thus its entire top-form space has dimension six. The refinement therefore applies to reduced boundaries; the
character-line construction of Section~\ref{rankone:section} is used
for general rational orbifold boundaries.
\end{remark}

\subsection{The original relative orders}


\begin{lemma}[Normalization of the rank-one boundary]
\label{lem:rankone-normalization}
Let $x:X'\to W$ be a projective fibration on smooth models, let $D'$ be
an effective reduced SNC divisor, and suppose that its geometric generic
fiber $(J,D_J)$ has logarithmic Kodaira dimension zero. Fix a nonzero
relative rational $p$-canonical form $s$ generating
$H^0(J,p(K_J+D_J))$. At a prime $P$ of $W$ set
\[
 a_E=\operatorname{ord}_E(x^*P),\qquad
 r_E=p^{-1}\operatorname{ord}^{pK_{X'/W}}_E(s),\qquad
 d_E=\operatorname{coeff}_E(D'),\qquad
 t_P=\min_{E\mapsto P}\frac{r_E+d_E}{a_E}.
\]
Here the minimum is over actual components on the chosen source model.
Choose that model so that the generator divisor, boundary and inverse
image of $P$ have SNC support over the generic point of $P$; later
source modifications carry strict boundary plus reduced exceptional
boundary. The coefficients in the rank-one comparison satisfy
\[
 \alpha_P=t_P,\qquad
 \beta_P=\inf_{E\mapsto P}\frac{1+r_E}{a_E}-1,\qquad
 T_P=1-c_P,
 \quad
 c_P:=\inf_{E\mapsto P}\frac{1+r_E}{a_E}-t_P.
\]
The infima allow divisorial valuations on higher smooth source models.
Equivalently, $c_P$ is the log canonical threshold of $x^*P$ for the
normalized auxiliary subpair
\[
 \Delta^{\mathrm{aux}}
 =-p^{-1}\operatorname{div}_{K_{X'/W}}(s)+x^*\sum_Qt_QQ.
\]
If $T_P=0$, an actual component attaining $t_P$ has
$a_E=1$, $d_E=0$, and attains the threshold infimum as well.
\end{lemma}

\begin{proof}
Choose a local frame $\eta$ of $K_W$ that does not vanish at the generic
point of $P$. The canonical line-bundle identification
\[
 \mathcal O_{X'}(pK_{X'/W})\otimes x^*\mathcal O_W(pK_W)
       \simeq\mathcal O_{X'}(pK_{X'})
\]
identifies the order of $s$ with the order of
$s\wedge x^*\eta^{\otimes p}$. Thus the quantity $l_E$ in the
rank-one order formula is exactly $r_E$. This uses the actual relative
canonical line bundle: the Jacobian contribution is already present in
$r_E$, so no ramification term is added or subtracted here.
Consequently its two order formulas give
\[
 \alpha_P=t_P,
 \qquad
 \beta_P=\min_E\frac{r_E+1-a_E}{a_E}.
\]
Put
\[
 \Delta_0=-p^{-1}\operatorname{div}_{K_{X'/W}}(s),
 \qquad \lambda_P=\min_{E\mapsto P}\frac{1+r_E}{a_E},
\]
where the minimum is over actual vertical components. The indicated
SNC preparation and the generic logarithmic bound give $\Delta_0$ SNC
support and horizontal coefficients at most one. In
$\Delta_0+\lambda_Px^*P$, each vertical coefficient is
$-r_E+\lambda_Pa_E\leq1$. Thus this subpair is sub-log-canonical
over the generic point of $P$, and every higher divisorial valuation
satisfies
\[
                     1+r_E-\lambda_Pa_E\geq0.
\]
An actual component attaining $\lambda_P$ gives equality. Hence the
infimum over all such valuations equals $\lambda_P$, and
$\beta_P=\lambda_P-1$.

For completeness, the first minimum also remains unchanged on higher
source models. A coefficient realizing $t_P$ makes
$t_Pa_E-r_E$ no greater than the original boundary coefficient at every
actual component over $P$. Equivalently, after clearing denominators,
the normalized generating form satisfies every logarithmic order test
over the generic point of $P$. Its pullback remains logarithmic with
strict plus reduced exceptional boundary. No new valuation can lower
the first minimum, while an old minimizing valuation persists.

At any divisorial valuation above the generic point of $P$, the
crepant coefficient of $\Delta^{\mathrm{aux}}$ is $-r_E+a_Et_P$.
The threshold formula therefore gives
\[
 \operatorname{lct}_{\eta_P}(X',\Delta^{\mathrm{aux}};x^*P)
 =\inf_E\frac{1+r_E-a_Et_P}{a_E}=c_P.
\]
Since $T_P=\alpha_P-\beta_P$, we obtain $T_P=1-c_P$.
Finally, if $T_P=0$ and $E$ attains the first minimum, then
\[
 1=c_P\le\frac{1+r_E}{a_E}-t_P
       =\frac{1-d_E}{a_E}\le1.
\]
All terms are equal. Since $a_E$ is a positive integer and
$d_E\in\{0,1\}$, this forces $a_E=1$ and $d_E=0$, as claimed.
\end{proof}

\begin{lemma}[Permitted models after normalization]
\label{lem:normalized-models}
Suppose that $(W,T)$ is smooth with effective SNC rational boundary and
that the moduli divisor $M$ is a rational Cartier divisor on $W$.
For a smooth modification $\mu:W'\to W$, put
\[
 T'=\mu_*^{-1}T+\operatorname{Exc}(\mu)_{\mathrm{red}},
 \qquad M'=\mu^*M,
\]
and assume the resulting boundary is SNC. Then pullback identifies the
section spaces of every sufficiently divisible multiple of
$K_W+T+M$ and $K_{W'}+T'+M'$. At unchanged codimension-one points the
normalization in Lemma~\ref{lem:rankone-normalization} is preserved.
Every new exceptional prime has coefficient one in $T'$.
\end{lemma}

\begin{proof}
Log canonicity of $(W,T)$ gives
\[
 K_{W'}+T'+M'=\mu^*(K_W+T+M)+E,
\]
where $E$ is effective and $\mu$-exceptional. For divisible $m$,
normality gives
\[
 \mu_*\mathcal O_{W'}(mE)=\mathcal O_W.
\]
The projection formula therefore identifies the section spaces.
The modification is an isomorphism at each old generic divisorial
point, so the old local order calculations are unchanged. The final
assertion is the definition of $T'$.
\end{proof}

Appendix~\ref{p1logcbfalt:section} gives complementary projective
proofs of the section comparisons and further calculations with the
auxiliary subpair on higher models.
